\section*{The Propers: Detailed Commentary}

\subsection*{Entrance: Doxology before Analysis\enspace\elementref{Entrance Antiphon}}

The Entrance Antiphon is a composed liturgical blessing rather than an attributed biblical quotation. It names Father, the Only Begotten Son, and Holy Spirit and gives one reason for praise: divine merciful love. Praise comes before analysis. The assembly does not begin by treating the mystery as an object placed at intellectual distance; it receives itself as a people already addressed by mercy and therefore capable of confession.

\subsection*{Collect: Missions Ordered to Profession\enspace\elementref{Collect}}

The Collect addresses the Father and recalls the sending of the Word of truth and the Spirit of sanctification. It then asks that servants profess the true faith by acknowledging the glory of the eternal Trinity and adoring Unity in almighty majesty. The prayer's order moves from divine action to ecclesial confession: revelation and sanctification precede profession.

\begin{threestable}{Prayer movement}{Doctrinal work}{Distinction}
Word of truth sent & The Son's mission is revelatory and personal, not the arrival of an impersonal message. & The mission in time does not begin the Son's existence.\\
Spirit of sanctification sent & Knowledge of God is joined to transformation and worship. & The Spirit is not reduced to an effect, mood, or created force.\\
Trinity acknowledged; Unity adored & Real distinction and one divine life are confessed together. & Neither three gods nor one Person wearing three names is admitted.\\
\end{threestable}

Catechism 236 draws the distinction: the economy reveals the divine life without exhausting it. The prayer can move from sending to confession because the temporal missions manifest who the Son and Spirit eternally are; it cannot license collapsing eternal processions into visible historical events.

\begin{sourcecard}{St. Gregory Nazianzen --- \emph{Oration 40}, 41--43}
In baptismal preaching Gregory moves contemplation from the One to the Three and back without allowing either movement to cancel the other. Unity is not a solitary Person; distinction is not partition of power or Godhead. His doxological rhythm illuminates the Collect's paired verbs “acknowledge” and “adore”: the mystery is confessed with precision because it is given for worship, not because a verbal formula encloses divine life.
\end{sourcecard}

Gregory's setting also keeps the prayer ecclesial. The confession arises around baptism and Christian initiation, not as an exercise for a detached specialist. Yet baptismal accessibility does not justify loose language. The newly initiated are invited into the Church's exact rule because Father, Son, and Spirit are not metaphors for three human experiences of God.

\subsection*{Exodus and Daniel: Mercy Answered by Blessing\enspace\elementref{First Reading, Responsorial Canticle}}

Exodus 34 comes after rupture. Israel has worshipped the calf; Moses has shattered the tablets; intercession and threatened absence expose the covenant's crisis. The appointed passage does not minimize sin. It places pardon within God's free self-proclamation and Moses' bold request that the Lord accompany, forgive, and take this people as an inheritance.

The divine proclamation belongs first to that covenant-renewal setting. Repeated divine names and attributes do not form a concealed numerical code for three Persons. Cyril of Jerusalem nevertheless provides exact Christian reception in \emph{Catechetical Lecture X} 7--8: he quotes the descent, proclamation of the divine name, mercy, Moses' bow, and the plea that the Lord go amid the people. Cyril uses the sequence while teaching about Father and Son. He does not count titles into Persons; he reads the Lord whose merciful presence Moses seeks within the Church's fuller Christological rule.

Gregory of Nyssa, in \emph{Against Eunomius} XI, likewise quotes the adjoining Moses narrative through Exodus 34:9 while arguing that the divine messenger who accompanies Israel cannot be reduced to a creature. His reading is a fourth-century Christological reuse governed by later controversy, rather than a continuous literal exposition. In both witnesses the Old Testament's one divine mercy enters Christian doctrinal argument without speaking technical Nicene prose in advance.

\begin{fourstable}{Textual moment}{Literal action}{Liturgical fruit}{Distinction}
Replacement tablets & Covenant words are received again after infidelity. & The assembly hears mercy as restoration, not denial of sin. & Grace is not permission to call the calf harmless.\\
Divine proclamation & God names mercy, fidelity, justice, and patience. & Trinitarian worship begins from the one God's revealed identity. & Repeated titles are not three hidden Persons.\\
Moses bows & Revelation elicits bodily worship and intercession. & Doctrine tends toward adoration and prayer for a people. & Mystery is not mastered as abstract data.\\
Daniel responds & The persecuted bless God across heaven, sanctuary, throne, and depths. & The response widens covenant mercy into doxology. & The canticle survives in Greek, outside the Hebrew-Aramaic Daniel.\\
\end{fourstable}

\Needspace{5\baselineskip}
The Daniel verses belong to a later Greek addition, yet the Church receives them canonically as Scripture. Their late literary formation does not weaken their liturgical authority, and their canonical authority does not depend on a Babylonian composition date. Their responsorial force is plain: before asking how the one God is three Persons, the assembly blesses the God whose mercy makes worship possible.

Cyril of Jerusalem directly cites the Song's language in \emph{Catechetical Lecture IX} 3. The God who beholds the depths and sits above the Cherubim is, for Cyril, the one Creator of visible and invisible reality; the verse excludes rival creators rather than supplying a Trinitarian count. This reception respects the canticle's doxological grammar: heaven, sanctuary, throne, firmament, and depths do not divide divine rule into regions.

Chrysostom also appeals to the furnace prayer in \emph{Homily 2 on Philemon}, emphasizing the three youths' humility in affliction. His citation falls elsewhere in the Song than the appointed 3:52--56 and is therefore contextual reception, not direct commentary. Together the witnesses show two ancient uses of the same Greek addition---cosmic confession and penitential humility---while the responsorial verses themselves remain a blessing of the one God.

The complete five-verse unit also has unusually dense liturgical reception as the \emph{Benedictus es} canticle across ancient and medieval Western offices and rites. Indexed sources preserve hundreds of uses of the sequence. Those use records are not commentary, and no secure continuous patristic exposition of all five verses has been located in the checked corpus; they show how generations learned the verses by repeated doxology. The canticle's heaven, glory, temple, throne, and firmament became a stable grammar of praise before their appointment in this modern responsorial setting.

Exegesis and liturgical reception differ here. Cyril supplies verse-level doctrine and the canticle's long ecclesial use supplies a practice of praise, but neither is a continuous interpretation by an unnamed “tradition.” The most certain traditional fact is the act of blessing itself: the Church repeatedly returns God's revealed names and cosmic rule as praise.

\subsection*{Second Reading: A Blessing with Ethical Syntax\enspace\elementref{Second Reading}}

Paul begins the pericope with imperatives and closes with greeting and blessing. Restoration, encouragement, agreement, and peace are not detachable preliminaries. They are the communal shape within which grace, divine love, and the Spirit's communion are invoked. The God of love and peace is promised in a sentence that also commands the wounded Corinthians to seek reconciliation.

\begin{sourcecard}{St. John Chrysostom --- \emph{Homily 30 on 2 Corinthians}, 3}
Commenting directly on the final blessing, Chrysostom observes the Father, Son, and Holy Spirit named together and refuses to partition their gifts as though each Person acted independently. His older numbering says 13:14; the U.S. Lectionary says 13:13.
\end{sourcecard}

The text is a strong apostolic ground for Trinitarian confession in the order Paul gives it. “Grace” does not make the Son lesser than the Father; “love” does not belong to the Father alone; “communion” does not reduce the Spirit to a shared human feeling. Augustine's rule of inseparable operation and the proper Preface's equality of majesty keep the invocation from being divided into three competing benefactors.

Chrysostom's close is also ecclesial rather than merely taxonomic. He reads the threefold gift after Paul's commands to restoration and peace, and he argues from the coordination of the divine Persons rather than assigning each a sealed department. The liturgical formula accordingly joins doctrine and repaired communion: a congregation cannot confess inseparable divine action while treating grace, love, and fellowship as excuses for division.

The greeting therefore functions as blessing and examination. Grace asks whether relationships are being restored by gift rather than domination; love asks whether agreement is truthful rather than coerced; communion asks whether peace has concrete ecclesial form. Those three questions come from Paul's imperative context; none of them assigns a modern program to a divine Person. The one divine work creates differentiated gifts that return the community toward one peace.

\subsection*{Acclamation and Gospel: The Saving Gift in the Light\enspace\elementref{Gospel acclamation, Gospel}}

The acclamation is explicitly marked \emph{cf.} Revelation 1:8. It takes the eternal divine title and gives it a Trinitarian liturgical form. The adaptation is legitimate appointed text, but the added formula is the Lectionary's, not what Revelation's author says in that verse. Acclamatory function, not proof-text arithmetic, governs its place.

Victorinus of Pettau's early Latin \emph{Commentary on the Apocalypse} receives the opening temporal title as a confession of the eternal one who creates with the Father and comes for judgment. He does not comment on the Lectionary's later Trinitarian adaptation, which did not yet exist. His witness therefore belongs to Revelation's source phrase only: it confirms that the title carried creation and judgment claims before the liturgy placed it in an explicitly triune doxology.

John 3:16--18 stands within a conversation about birth from above, Spirit, testimony, the heavenly descent of the Son of Man, and the lifted serpent. The appointed verses concentrate on divine initiative: God loves, gives, and sends; the Son is received in faith; eternal life is offered. The exclusion of verses 19--21 does not erase the light-and-darkness explanation that immediately follows.

\begin{sourcecard}{St. Augustine --- \emph{Tractate 12 on John}, 11--13}
Augustine reads the lifted serpent as a healing sign and insists that Christ's coming is ordered to salvation. Condemnation is disclosed where the sick person refuses the offered remedy and avoids the light. His reading holds together divine mercy, the Cross, faith, confession, and moral conversion.
\end{sourcecard}

\begin{sourcecard}{St. John Chrysostom --- \emph{Homilies 27--28 on John}}
Chrysostom directly follows John 3:16--18. The distance between the immortal giver and an ungrateful world magnifies love; “gave” points toward the Son's saving self-offering. The next homily distinguishes the purpose of the first coming from the judgment disclosed by unbelief: rescue is genuinely offered, yet refusal is not neutral.
\end{sourcecard}

Cyril of Jerusalem uses John 3:16 and 18 in \emph{Catechetical Lecture XI} 6 to confess the Son as only-begotten by nature rather than an adopted creature. Cyril's catechesis concerns the Son's nature; John's own argument concerns saving love. Cyril's reading clarifies why the feast may hear the Father's gift as revelation of who the Son eternally is while still letting the Gospel's immediate movement run from love through sending, faith, salvation, and judgment.

The two great Johannine preachers place different accents. Augustine organizes the verses around remedy and light; Chrysostom around the immeasurable condescension of giver and gift; Cyril around the natural sonship of the Only Begotten. None turns “world” into a claim that every response is identical, and none makes judgment the Father's motive for sending. Read together, they widen the passage without flattening it.

\begin{studybox}{John 3:16--18 and the feast}
The appointed Gospel directly names God and the only Son; it does not name all three Persons within verses 16--18. The Spirit belongs to the preceding discourse and the whole Johannine witness. The feast's Trinitarian synthesis comes from the canon and the liturgy, not from the pericope's isolated wording. Likewise, universal saving purpose does not imply that faith and judgment are unreal.
\end{studybox}

\subsection*{Prayer over the Offerings: The Gift Returned\enspace\elementref{Prayer over the Offerings}}

The proper prayer asks sanctification through invocation of the divine Name and moves from the gifts to the worshippers themselves as an enduring offering. It does not make human self-dedication an independent saving sacrifice. The Eucharistic action precedes and forms ecclesial self-offering; worshippers ask God to accomplish what they cannot manufacture.

\begin{sourcecard}{St. Augustine --- \emph{City of God} X.5--6, 20}
Augustine is not commenting on this prayer. He supplies doctrinal illumination: visible sacrifice signifies the interior offering that joins a person to God, and the redeemed city offers itself through Christ the one High Priest and sacrifice. The order preserves both Christ's uniqueness and the Church's real participation.
\end{sourcecard}

\subsection*{Proper Preface: The Rule that Guards the Feast\enspace\elementref{Preface}}

The Preface prevents the formula's biblical variety from becoming a loose collection of devotional images. It confesses one God and Lord, not one solitary Person; distinction belongs to Persons, unity to substance, equality to majesty. This is liturgical dogma within thanksgiving, not a speculative appendix after the practical texts.

\begin{fourstable}{Confessed term}{Affirmation}{Error excluded}{Effect on reading}
Three Persons & Father, Son, and Spirit are really distinct. & Modalism or three passing roles. & Missions and triadic blessings concern personal subjects.\\
One substance & The divine nature is numerically one. & Tritheism or a society of gods. & Mercy and saving operation are one divine work.\\
Equal majesty & No Person is lesser in divinity. & Subordinationism. & Sending and reception do not imply created inferiority.\\
One adoration & The Church's thanksgiving is directed to the triune God. & Partitioned worship. & Biblical strands end in one doxology.\\
\end{fourstable}

Augustine's \emph{De Trinitate} I.4.7 gives the Latin doctrine behind these confessed terms. Father, Son, and Spirit are one substance in indivisible equality; only the Son becomes incarnate, yet the external divine work is inseparable. The Father can therefore appropriately be said to give and send, the Son to be sent and incarnate, and the Spirit to sanctify, without three agencies negotiating a common project (Leo XIII, \emph{Divinum illud munus} 3).

Gregory Nazianzen supplies a complementary Greek accent. His return from unity to three and from three to unity is not numerical oscillation but protection against opposite distortions. Starting from one may hide the Persons; starting from three may divide the Godhead. The Preface's thanksgiving performs the same discipline by coordinating personal distinction, one substance, equal majesty, and one adoration in a single act.

This doctrine does not erase the biblical idioms. Exodus can still speak covenantal mercy, John of the Son given, Paul of grace, love, and communion, and Galatians of sending and filial cry. Dogmatic grammar orders their canonical hearing; it does not replace each text with a technical paraphrase. Conversely, appeal to “biblical language” cannot license a reading that contradicts the one Godhead or the real distinctions the Church confesses.

\Needspace{7\baselineskip}
\subsection*{Communion: Adoption as Trinitarian Life\enspace\elementref{Communion Antiphon}}

Galatians 4:6 can be misheard as a decorative proof text. Paul instead places divine action before the heir's response: the Son is sent for redemption and adoption, the Spirit of the Son into hearts, and filial address to the Father follows. Communion thus becomes relational without becoming private intimacy.

Chrysostom's \emph{Homily 4 on Galatians} directly follows 4:1--7. The Spirit's filial cry is evidence that those once enslaved have truly received adoption and become heirs; returning to bondage would contradict a gift already enacted. His emphasis prevents “Abba” from becoming an isolated emotion. The cry belongs to an objective change of status and a new manner of life.

Basil's \emph{On the Holy Spirit} uses Galatians 4:6 to explain participation in the Spirit: Scripture can speak of the Spirit crying because the Spirit makes the recipients' own voice capable of filial address. Augustine similarly says in \emph{Sermon 21 on the New Testament} 29 that the Spirit makes believers cry by pouring divine love into their hearts. These Greek and Latin witnesses avoid two opposite reductions. The Spirit is neither a created voice replacing the human person nor merely the person's religious self-expression; divine agency brings the adopted person's own prayer to life.

\begin{threestable}{Pauline movement}{Trinitarian form}{Moral consequence}
The Son is sent & Redemption and adoption are received through Christ. & Christian freedom is not self-originating autonomy.\\
The Spirit is sent into hearts & Divine communion becomes interior and ecclesial life. & The Spirit is more than a religious emotion.\\
The Father is addressed & The baptized live as heirs rather than slaves. & Filial confidence must become obedience and peace.\\
\end{threestable}

\subsection*{Prayer after Communion: Confession for the Whole Person\enspace\elementref{Prayer after Communion}}

The final prayer joins two realities without confusing them: reception of the sacrament and confession of the eternal Trinity and undivided Unity. Both are asked to advance salvation of body and soul. The prayer therefore refuses a disembodied intellectualism and a sacramental automatism. Reception and confessed faith are gifts ordered toward continuing transformation of the whole person.

“Body and soul” is not a division into a disposable shell and a religious core. The Gospel promises eternal life to embodied persons; Communion voices adoption within the Spirit; Paul asks a conflicted community to repair its common life. The prayer gathers intellectual confession, sacramental reception, bodily conduct, affect, and relationship into one saving trajectory. Its request remains necessary after Communion because a true gift can be resisted, incompletely integrated, or contradicted in practice.
